\documentclass[11pt]{article}
\usepackage[margin=1in]{geometry}
\usepackage[utf8]{inputenc}
\usepackage[T1]{fontenc}
\usepackage{amsmath,amssymb,amsfonts,amsthm,mathtools}
\usepackage{booktabs}
\usepackage{microtype}
\usepackage{enumitem}
\usepackage{xcolor}
\usepackage{tikz}
\usepackage[colorlinks=true,linkcolor=blue!60!black,citecolor=blue!60!black,urlcolor=blue!60!black]{hyperref}

% ---- theorem-like environments ----
\theoremstyle{plain}
\newtheorem{theorem}{Theorem}[section]
\newtheorem{lemma}[theorem]{Lemma}
\newtheorem{corollary}[theorem]{Corollary}
\newtheorem{proposition}[theorem]{Proposition}
\newtheorem{conjecture}[theorem]{Conjecture}
\theoremstyle{definition}
\newtheorem{definition}[theorem]{Definition}
\newtheorem{example}[theorem]{Worked example}
\newtheorem{problem}{Problem}[section]
\theoremstyle{remark}
\newtheorem{remark}[theorem]{Remark}

% ---- problem / hint / solution / discovery styling (printable, no external boxes) ----
\definecolor{hintc}{rgb}{0.35,0.35,0.35}
\definecolor{movec}{rgb}{0.0,0.30,0.55}
\definecolor{discc}{rgb}{0.55,0.20,0.0}
\newenvironment{hints}{\par\smallskip\noindent\textbf{\color{hintc}Hints (peek one at a time):}\begin{itemize}}{\end{itemize}\par\smallskip}
\newenvironment{solution}{\par\smallskip\noindent\textit{Solution.}\;}{\hfill$\square$\par\smallskip}
\newcommand{\researchmove}[2]{\par\smallskip\noindent{\color{movec}\rule{\linewidth}{0.4pt}}\par\nobreak
  \noindent{\color{movec}\textbf{Research move: #1}}\par\nobreak\noindent #2\par
  \noindent{\color{movec}\rule{\linewidth}{0.4pt}}\par\smallskip}
\newcommand{\discover}[2]{\par\smallskip\noindent{\color{discc}\textbf{Find-the-question: #1}}\;#2\par\smallskip}

% ---- shared macros (match the technical report) ----
\newcommand{\R}{\mathbb{R}}
\newcommand{\Z}{\mathbb{Z}}
\newcommand{\E}{\mathbb{E}}
\newcommand{\fdc}{f_{\mathrm{dc}}}
\newcommand{\Vinv}{V_{\mathrm{inv}}}
\newcommand{\Vac}{V_{\mathrm{ac}}}
\newcommand{\Sep}{\mathrm{Sep}}
\newcommand{\one}{\mathbf 1}
\newcommand{\Piinv}{\Pi_{\mathrm{inv}}}
\newcommand{\sign}{\mathrm{sign}}
\newcommand{\conv}{\operatorname{conv}}
\newcommand{\dist}{\operatorname{dist}}
\newcommand{\Lip}{\operatorname{Lip}}
\newcommand{\Orb}{\mathcal O}

\title{\textbf{Shift-Invariance and Adversarial Robustness}\\[3pt]
\large A Collaborator Primer: background, problem-style proofs, and how to find the next question\\[4pt]
\normalsize companion to the paper and the technical report; for a new theory collaborator}
\author{Anass Al Ammiri}
\date{\today}

\begin{document}
\maketitle

\begin{center}\small\itshape
A guided path that takes a strong undergraduate or new graduate student from zero to (i) reading the paper and
the technical report end to end, (ii) reconstructing every proof as a sequence of problems they solve themselves,
and (iii) learning the harder skill: how to find the questions, not only follow the steps. Plan three to four weeks
of focused study. Companion to \texttt{paper/report/main.tex} (the paper) and \texttt{paper/report/theory\_full\_checked\_v2.tex}
(the technical report, every proof in full).
\end{center}

\medskip\hrule\medskip

\section*{Preface}
\addcontentsline{toc}{section}{Preface}

\paragraph{Who this is for.} A new collaborator joining the shift-invariance/adversarial-robustness project,
mostly on the theory side and a little on the experiments. We assume one undergraduate course each in linear
algebra, real analysis or advanced calculus, and probability, plus some exposure to machine learning. You do
\emph{not} need prior background in adversarial robustness, Fourier analysis on groups, or implicit-bias theory;
all three are built from first principles in Part~I.

\paragraph{What this is, and what makes it different.} Most primers decompose each proof into steps and ask you to
follow them. Following steps is the easy half of mathematics. The hard half, the half that actually produces
research, is \emph{coming up with the steps yourself}: knowing which question to ask, which object to define,
which special case to test first. This document trains both. Every result appears first as a \textbf{Problem} you
attempt cold, then a graded ladder of \textbf{Hints} you uncover only when stuck, then a full \textbf{Solution}.
Interleaved are two recurring devices:
\begin{itemize}[leftmargin=1.4em,nosep]
\item \textbf{Research move} boxes that name the general technique a proof secretly used (``quotient out the
nuisance,'' ``diagonalize the symmetry,'' ``find the degenerate case''), so you can reuse it.
\item \textbf{Find-the-question} prompts that hand you a situation with \emph{no question attached} and ask you to
generate the right one before reading on. This is the skill that does not come from following steps.
\end{itemize}

\paragraph{How to use it.} Read Part~I in order; do the problems before reading their solutions, and give each at
least fifteen honest minutes before the first hint. Part~II is the paper's theory, result by result; by the end
you should be able to state and prove each result with the book closed. Part~III is the discovery layer; do it
\emph{after} Part~II, because it reflects on how the Part~II results were found. Part~IV points at the open
problems where you can contribute and how to run the experiments.

\paragraph{House writing rules (we follow them, so should you).} Plain language. Define every term on first use.
No em-dashes, no ``win''/``beat'' as verbs, no AI-tell filler (``delve,'' ``leverage,'' ``crucial''), one spelling
convention per document, captions that match figures, honest scoped claims. A claim you cannot verify against the
paper, the technical report, or a results file is a claim to verify, not to repeat.

\paragraph{One running thread.} A single object recurs from Part~I to Part~III: the cyclic group of circular shifts
acting on a length-$d$ signal, its \emph{orbit} (the set of all shifts of one signal), and the split of $\R^d$ into
the constant ``DC'' line and its zero-sum complement. If you understand that one picture deeply, everything else is
commentary on it.

\newpage
\tableofcontents
\newpage

% =====================================================================
\section{Master notation table}
The symbols below are used throughout. Section references point to where each is introduced.

\begin{center}\small
\begin{tabular}{ll}
\toprule
Symbol & Meaning \\
\midrule
$x\in\R^d$ & input signal or flattened image; $d$ coordinates (\S\ref{sec:rob}) \\
$y\in\{\pm1\}$ & binary label; multiclass uses the active margin $M$ (\S\ref{sec:rob}) \\
$f:\R^d\to\R$ & score; prediction $\sign f(x)$ (\S\ref{sec:rob}) \\
$M(x)=F_y(x)-\max_{c\ne y}F_c(x)$ & active logit margin (positive iff correct) (\S\ref{sec:rob}) \\
$B_p(\varepsilon)$, $r_p(f,x)$ & threat ball of radius $\varepsilon$; robust radius (distance to boundary) (\S\ref{sec:rob}) \\
$L$, $\alpha$, $\kappa=L/\alpha$ & upper-Lipschitz, co-Lipschitz constants; their ratio (\S\ref{sec:rob}, \S\ref{sec:etaL}) \\
$\eta$ & feature-space margin of an invariant feature (\S\ref{sec:etaL}) \\
$\eta/L$ & margin-to-Lipschitz ratio; the certificate scalar (\S\ref{sec:etaL}) \\
$G=\Z_d$, $S_s$ & cyclic group of circular shifts; the shift-by-$s$ operator (\S\ref{sec:group}) \\
$\Orb(x)=\{S_sx\}$ & orbit of $x$ (all its shifts) (\S\ref{sec:group}) \\
$\hat x_k$, $|\hat x_k|^2$ & unitary DFT coefficient; power spectrum (\S\ref{sec:group}) \\
$B_{k_1,k_2}(x)$ & bispectrum, a degree-three shift invariant (\S\ref{sec:group}) \\
$\Vinv=\mathrm{span}(\one)$, $\Vac$ & DC (constant) line and its zero-sum complement (\S\ref{sec:group}) \\
$\Piinv=ee^\top$, $e=d^{-1/2}\one$ & projector onto the DC line (\S\ref{sec:group}) \\
$P_G=|G|^{-1}\sum_g R_g$ & group-average (Reynolds) projector onto the fixed subspace (\S\ref{sec:group}) \\
$\fdc(x)=d^{-1/2}\sum_i x_i$ & unit DC functional (\S\ref{sec:group}) \\
$\gamma_{\mathrm{inv}}$, $\Sep$ & invariant linear margin; separation of two sets (\S\ref{sec:margin}, \S\ref{sec:thmA}) \\
$Q(x)=(|\hat x_k|^2)_k$ & power-spectrum feature map (\S\ref{sec:ps}) \\
$f_\theta(x)=\sum_r a_r\sigma(w_r^\top x)$ & two-layer ReLU network, degree-2 homogeneous (\S\ref{sec:bias}) \\
$\mathcal H_K$, $\Pi_G$ & RKHS of the tangent kernel; its orbit-averaging projection (\S\ref{sec:lazy}) \\
\bottomrule
\end{tabular}
\end{center}

% =====================================================================
\part*{Part I: Background you need before the paper}
\addcontentsline{toc}{section}{Part I: Background you need before the paper}

This part assembles the three toolkits the paper rests on: adversarial robustness, Fourier analysis on the cyclic
group, and margin/implicit-bias theory. Each subsection states the primitives, works a numerical example you can
check by hand, and ends with problems. Skim a subsection if you already know it, but do its problems; they are
chosen to expose the exact facts Part~II reuses.

\section{Adversarial robustness from scratch}\label{sec:rob}

\paragraph{The object of study.} A classifier is a score $f:\R^d\to\R$ predicting $\sign f(x)$ (multiclass: a vector
$F$ predicting $\arg\max_c F_c$). Standard training maximizes clean accuracy. The robustness question asks something
stricter: if an adversary may move $x$ by a small vector $\delta$, can it flip the prediction?

\begin{definition}[Threat model, robust radius, robust accuracy]
An $\ell_p$ threat of budget $\varepsilon$ allows $\delta\in B_p(\varepsilon)=\{\delta:\|\delta\|_p\le\varepsilon\}$.
The two standard norms are $\ell_2$ (Euclidean) and $\ell_\infty$ (each pixel moves by at most $\varepsilon$). The
\emph{robust radius} at a correctly classified $(x,y)$ is $r_p(f,x)=\inf\{\|\delta\|_p:\sign f(x+\delta)\ne y\}$,
the distance to the decision boundary. The \emph{robust accuracy at $\varepsilon$} is the fraction of test points
with $r_p>\varepsilon$.
\end{definition}

The radius is the finer quantity; robust accuracy is the radius thresholded at $\varepsilon$, and it shows floor
effects (everything reads zero) when $\varepsilon$ is large relative to typical radii. The paper uses the radius
as the dataset-comparable measure and robust accuracy only where a standard threshold makes it meaningful (after
adversarial training).

\paragraph{Attacks versus certificates.} An \emph{attack} (PGD, AutoAttack, the minimum-norm DDN attack) searches
for a small $\delta$ and so \emph{upper bounds} the radius: it can only show the model is at least this weak. A
\emph{certificate} \emph{lower bounds} the radius: it guarantees the model is at least this strong. The two bracket
the truth from opposite sides. The single most important certificate is the next definition.

\begin{definition}[Lipschitz--margin certificate]\label{def:lipmargin}
A score $g$ is \emph{$L$-Lipschitz} toward the boundary if $|g(x)-g(z)|\le L\|x-z\|$ along the relevant paths. If
$g$ is correct at $x$ with score margin $\mu=y\,g(x)>0$, then no $\delta$ with $\|\delta\|<\mu/L$ can drive $g$ to
zero, so $r(g,x)\ge\mu/L$. For a differentiable score the first-order $L$ is the dual-norm gradient
$\|\nabla_x M(x)\|_q$ ($q$ dual to the threat $p$: $\ell_\infty$ pairs with $\ell_1$, $\ell_2$ with $\ell_2$).
\end{definition}

\begin{example}[A one-line certificate]
Let $g(x)=w^\top x$ be linear, correct at $x$ with $\mu=y\,w^\top x$. Then $g$ is exactly $\|w\|_2$-Lipschitz in
$\ell_2$, so $r_2\ge\mu/\|w\|_2$. In fact for a linear score this is an equality: the distance from a point to a
hyperplane is exactly $\mu/\|w\|_2$. Hold onto that ``equality for the linear case''; \S\ref{sec:etaL} shows it is
the boundary case of a more general two-sided bound.
\end{example}

\begin{problem}[Why the dual norm]\label{pr:dual}
Show that for a differentiable score the largest first-order change $|g(x+\delta)-g(x)|$ over $\|\delta\|_p\le\varepsilon$
is $\varepsilon\,\|\nabla g(x)\|_q$ with $1/p+1/q=1$, and deduce the first-order radius $M(x)/\|\nabla M(x)\|_q$.
\end{problem}
\begin{hints}
\item First order, $g(x+\delta)-g(x)\approx \nabla g(x)^\top\delta$.
\item Maximizing a linear functional over an $\ell_p$ ball is the definition of the dual ($\ell_q$) norm.
\item For $\ell_\infty$ the maximizer puts $\delta_i=\varepsilon\,\sign(\partial_i g)$, giving $\varepsilon\|\nabla g\|_1$.
\end{hints}
\begin{solution}
By Cauchy--Schwarz's $\ell_p$/$\ell_q$ generalization (Hölder), $\sup_{\|\delta\|_p\le\varepsilon}\nabla g^\top\delta
=\varepsilon\|\nabla g\|_q$, attained at the aligning $\delta$. The prediction flips when the change cancels the
margin, i.e.\ at $\varepsilon=M(x)/\|\nabla M(x)\|_q$ to first order. The pairing $\ell_\infty\!\leftrightarrow\!\ell_1$
is why an $\ell_\infty$ attack is read with the $\ell_1$ gradient norm, a fact the adversarial-training experiments
lean on heavily.
\end{solution}

\researchmove{Match the dual to the threat}{A robustness statement is meaningless until the norm is fixed, and the
right companion of an $\ell_p$ attack is the $\ell_q$ gradient. Whenever you see $\|\nabla M\|$ in this project, ask
``in which norm, and dual to which threat?'' The paper's adversarial-training result lives or dies on this matching:
the $\ell_1$ gradient predicts $\ell_\infty$ robustness, the $\ell_2$ gradient predicts $\ell_2$ robustness.}

\section{Group actions, Fourier analysis, and invariance}\label{sec:group}

\paragraph{Shifts as a group.} The cyclic group $G=\Z_d$ acts on $\R^d$ by circular shifts $(S_sx)_i=x_{(i-s)\bmod d}$.
Each $S_s$ is an orthogonal permutation (it just renumbers coordinates). The \emph{orbit} of $x$ is
$\Orb(x)=\{S_sx:s\in\Z_d\}$, the set of all $d$ shifts of $x$. A set is \emph{orbit-closed} if it contains the full
orbit of each of its members; a function is \emph{$G$-invariant} if $f(S_sx)=f(x)$ for all $s$.

\paragraph{Why Fourier.} All shift operators $S_s$ commute and are simultaneously diagonalized by the discrete
Fourier transform (DFT). With the unitary convention $\hat x_k=d^{-1/2}\sum_j x_j e^{-2\pi ijk/d}$, a shift
multiplies coefficient $k$ by the unit-modulus phase $e^{-2\pi isk/d}$. Two consequences run through the whole paper:
\begin{itemize}[leftmargin=1.4em,nosep]
\item The \emph{power spectrum} $|\hat x_k|^2$ is shift-invariant (the phase cancels in modulus) but \emph{phase-blind}.
\item The \emph{bispectrum} $B_{k_1,k_2}(x)=\hat x_{k_1}\hat x_{k_2}\overline{\hat x_{k_1+k_2}}$ is also shift-invariant
(the three shift phases cancel because $k_1+k_2-(k_1+k_2)=0$) and, unlike the power spectrum, keeps phase, so it
separates signals the power spectrum confuses.
\end{itemize}

\begin{definition}[Fixed subspace and the group-average projector]
The fixed subspace is $V^G=\{v:R_gv=v\ \forall g\}$ and the \emph{group-average (Reynolds) projector} is
$P_G=|G|^{-1}\sum_g R_g$, the orthogonal projector onto $V^G$. For cyclic shifts the only vectors fixed by every
shift are the constants, so $V^G=\Vinv=\mathrm{span}(\one)$ is the ``DC line,'' $P_G=\Piinv=ee^\top$ with
$e=d^{-1/2}\one$, and $\Vac=\Vinv^\perp=\{v:\sum_i v_i=0\}$ is the zero-sum ``AC'' subspace.
\end{definition}

\begin{example}[Hand-check at $d=4$]
Take $x=(1,0,0,0)$. Its orbit is the four standard basis vectors. $\fdc(x)=\tfrac12\sum_i x_i=\tfrac12$ for every
orbit member, so the DC coordinate is constant on the orbit (it must be: it is shift-invariant). The power spectrum
of $x$ is $|\hat x_k|^2=1/4$ for all $k$ (a spike is flat in frequency), so the power spectrum cannot tell $x$ from
any other unit spike, and in particular not from $-x$ shifted. This single fact is why the localized ``dot'' example
of the base paper is the hardest case for invariance.
\end{example}

\begin{problem}[The projector is the average and it cannot lengthen a vector]\label{pr:proj}
Prove that $P_G=|G|^{-1}\sum_g R_g$ is (i) idempotent, (ii) self-adjoint, hence an orthogonal projector, and (iii)
non-expansive: $\|P_Gv\|\le\|v\|$. Then specialize to cyclic shifts to show $P_G=ee^\top$.
\end{problem}
\begin{hints}
\item For (i) use the group property: $R_gP_G=P_G$ because $g\mapsto gh$ permutes the group.
\item For (ii) use that each $R_g$ is orthogonal, so $R_g^\top=R_{g^{-1}}$, and summing over the group is symmetric.
\item An orthogonal projector never increases the norm by Pythagoras: $\|v\|^2=\|P_Gv\|^2+\|(I-P_G)v\|^2$.
\item For cyclic shifts, $\frac1d\sum_s S_s$ sends every coordinate to the average, i.e.\ the matrix $\frac1d\one\one^\top=ee^\top$.
\end{hints}
\begin{solution}
(i) $P_G^2=|G|^{-2}\sum_{g,h}R_gR_h=|G|^{-2}\sum_{g,h}R_{gh}=|G|^{-1}\sum_k R_k=P_G$, since for fixed $g$, $h\mapsto gh$
is a bijection of $G$. (ii) $R_g^\top=R_g^{-1}=R_{g^{-1}}$, so $P_G^\top=|G|^{-1}\sum_g R_{g^{-1}}=P_G$. A self-adjoint
idempotent is an orthogonal projector; its range is exactly $V^G$ because $P_Gv=v\Leftrightarrow R_gv=v\ \forall g$.
(iii) Pythagoras gives $\|v\|^2=\|P_Gv\|^2+\|(I-P_G)v\|^2\ge\|P_Gv\|^2$. For cyclic shifts $\frac1d\sum_s S_s$ has
every entry $1/d$, i.e.\ $ee^\top$, projecting onto $\mathrm{span}(\one)$.
\end{solution}

\researchmove{Diagonalize the symmetry}{Whenever one symmetry group acts, look for the basis that diagonalizes it
all at once. For commuting shifts that basis is the Fourier basis, and ``invariant'' becomes ``depends only on
magnitudes,'' ``equivariant'' becomes ``multiply each frequency by a phase.'' Half the theorems in this project are
one line once you are in Fourier coordinates and three pages otherwise. The non-expansiveness in
Problem~\ref{pr:proj}(iii) returns, unchanged, as the engine of the lazy-regime theorem in \S\ref{sec:lazy}.}

\discover{a phase-blind invariant}{You are told the power spectrum is shift-invariant but cannot separate a spike
from its reflection. Before reading \S\ref{sec:ps}, ask: \emph{is there a shift-invariant feature that does keep
phase?} What is the lowest-degree polynomial in the $\hat x_k$ whose phase factors cancel under a shift? Write the
condition $k_1+k_2-k_3=0$ and see what degree-three monomial it forces. You will have re-derived the bispectrum.}

\section{Margins, convex hulls, and the max-margin classifier}\label{sec:margin}

\paragraph{Geometric margin.} For a linear score $w^\top x+b$, the geometric margin at $x$ is
$y(w^\top x+b)/\|w\|_2$, the signed Euclidean distance from $x$ to the hyperplane. The hard-margin classifier of two
finite sets maximizes the smallest margin.

\begin{theorem}[Convex-hull duality, classical]\label{thm:hull}
The maximum margin separating two finite sets $X_+,X_-$ equals $\tfrac12\dist(\conv X_+,\conv X_-)$, half the
distance between their convex hulls, attained by the hyperplane that perpendicularly bisects the closest pair of hull
points.
\end{theorem}

We will not reprove this classical fact, but we will use it constantly: it turns ``what is the best linear
separator'' into ``how far apart are the two clouds, after whatever transformation we applied.'' Every margin
computation in Part~II is an application of Theorem~\ref{thm:hull} to a transformed dataset (projected onto the DC
line, or pushed through the power-spectrum map).

\begin{problem}[One-dimensional margins]\label{pr:1d}
Two sets of real numbers, $X_+\subset[a_+,b_+]$ and $X_-\subset[a_-,b_-]$ with $b_-<a_+$. Show the max-margin
threshold sits at $(a_++b_-)/2$ with margin $\tfrac12(a_+-b_-)$, i.e.\ half the gap between the nearer endpoints.
\end{problem}
\begin{hints}
\item On the line, convex hulls are intervals; their distance is the gap between nearer endpoints.
\item Apply Theorem~\ref{thm:hull} with $\dist([a_-,b_-],[a_+,b_+])=a_+-b_-$.
\end{hints}
\begin{solution}
The hulls are $[a_-,b_-]$ and $[a_+,b_+]$; since $b_-<a_+$ they are disjoint with distance $a_+-b_-$, realized by the
endpoints $b_-$ and $a_+$. Theorem~\ref{thm:hull} gives margin $\tfrac12(a_+-b_-)$ and threshold at the midpoint
$(a_++b_-)/2$. This is the exact computation behind the invariant linear margin in \S\ref{sec:thmA}, where the line
is the DC axis.
\end{solution}

\section{Implicit bias: homogeneous nets, NTK, rich versus lazy}\label{sec:bias}

\paragraph{The selection problem.} On separable data, infinitely many parameter settings fit perfectly. Which one
does gradient descent pick? For a linear model on separable data it converges in direction to the max-margin
separator (Soudry et al.). For a positively homogeneous network ($f_{c\theta}=c^k f_\theta$; a bias-free two-layer
ReLU net is degree-2), normalized gradient flow converges in direction to a Karush--Kuhn--Tucker (KKT) point of the
margin program $\min\tfrac12\|\theta\|^2$ s.t.\ $y_if_\theta(x_i)\ge1$ (Lyu--Li, Ji--Telgarsky). These results say
``some KKT point''; \emph{which} one when several exist is the open content the project's optimization section is about.

\paragraph{Two training regimes.} A very wide network in the \emph{lazy} (neural-tangent-kernel, NTK) regime stays
close to its linearization at initialization and behaves like a fixed kernel machine (Jacot et al.; Chizat et al.).
In the \emph{rich} (small-initialization) regime it learns features. The lazy regime is where we can prove things
cleanly (\S\ref{sec:lazy}); the rich regime is where the hardest open question lives.

\paragraph{Feature averaging.} On data with many nearly orthogonal cluster centers, the gradient-descent max-margin
solution aligns with a signed average of those centers (Frei et al.; Li et al.). That classifies cleanly but is
fragile along the individual feature directions, giving a small robust radius. Keep this picture: it is the
mechanism the project must either invoke or rule out on orbit-structured data.

\begin{problem}[Homogeneity and the min-norm representation of one ReLU unit]\label{pr:relu}
A single ReLU unit computes $a\,\sigma(w^\top x)$ with $\sigma(t)=\max(t,0)$. Show the function is unchanged under
$(a,w)\mapsto(a/c,\,cw)$ for $c>0$, and that the representation cost $\tfrac12(a^2+\|w\|^2)$ is minimized at
$c^2=|a|/\|w\|$, with minimum value $|a|\,\|w\|$.
\end{problem}
\begin{hints}
\item $\sigma$ is positively homogeneous of degree one: $\sigma(c\,t)=c\,\sigma(t)$ for $c>0$.
\item So $a\,\sigma(w^\top x)=(a/c)\,\sigma((cw)^\top x)$. The function is invariant; only the cost changes.
\item Minimize $\tfrac12(a^2/c^2+c^2\|w\|^2)$ over $c$; set the derivative to zero.
\end{hints}
\begin{solution}
By degree-one homogeneity $(a/c)\sigma((cw)^\top x)=(a/c)\,c\,\sigma(w^\top x)=a\sigma(w^\top x)$, so the function is
unchanged. The cost $\tfrac12(a^2/c^2+c^2\|w\|^2)$ is convex in $c^2$ with minimum where
$a^2/c^4=\|w\|^2$, i.e.\ $c^2=|a|/\|w\|$, giving value $\tfrac12\cdot2|a|\|w\|=|a|\,\|w\|$. This ``$|a|\,\|w\|$ per
unit'' bookkeeping is the entire engine of the self-symmetrization theorem in \S\ref{sec:selfsym}; the proof there is
this problem applied to a whole orbit of tied units.
\end{solution}

% =====================================================================
\part*{Part II: The paper's theory, result by result}
\addcontentsline{toc}{section}{Part II: The paper's theory, result by result}

Each section states a result, hands it to you as a problem with a hint ladder, then gives the full solution. Do the
problem first. The technical report \texttt{theory\_full\_checked\_v2.tex} has the same proofs in reference form; use
it to check, not to substitute for the attempt.

\section{The invariant linear margin equals a projection separation}\label{sec:thmA}

The first structural result answers ``how much separation survives if the classifier must be shift-invariant?''

\begin{theorem}[Invariant linear margin]\label{thm:A}
A linear score $g(x)=w^\top x+b$ is shift-invariant for all $x$ iff $w\in\Vinv=\mathrm{span}(\one)$. Consequently,
for any finite separable classes $X_+,X_-$, the invariant linear max-margin equals
$\gamma_{\mathrm{inv}}=\tfrac12\dist(I_+,I_-)=\tfrac12\dist(\conv\Piinv X_+,\conv\Piinv X_-)$, where
$I_\pm=[\min_{x\in X_\pm}\fdc(x),\,\max_{x\in X_\pm}\fdc(x)]$ are the intervals of DC values.
\end{theorem}

\begin{problem}[Prove Theorem~\ref{thm:A}]\label{pr:A}
Two parts. (a) Characterize which $w$ give a shift-invariant linear score. (b) Reduce the invariant max-margin to a
one-dimensional problem and evaluate it.
\end{problem}
\begin{hints}
\item (a) Invariance for all $x$ means $w^\top S_sx=w^\top x$ for all $x,s$; turn ``for all $x$'' into a condition on $w$.
\item (a) $w^\top S_s x=(S_s^\top w)^\top x$, so you need $S_s^\top w=w$ for all $s$. Which vectors are fixed by every shift?
\item (b) With $w=c\one$, the score depends on $x$ only through $\fdc(x)$. The data collapse to scalars.
\item (b) Now it is Problem~\ref{pr:1d} on the DC axis.
\end{hints}
\begin{solution}
(a) Invariance $w^\top S_sx=w^\top x$ for all $x$ is equivalent to $(S_s^\top w-w)^\top x=0$ for all $x$, hence
$S_s^\top w=w$ for every $s$; the only vectors fixed by all cyclic shifts are constants, so $w\in\mathrm{span}(\one)$.
(b) Writing $w=c\one$, the score is a function of $\fdc(x)=\one^\top x/\sqrt d$ alone, so the classes reduce to the
scalar sets $\fdc(X_\pm)$, whose hulls are the intervals $I_\pm$. By Problem~\ref{pr:1d} the max-margin is half the
gap between the nearer endpoints, $\tfrac12\dist(I_+,I_-)$. Identifying the projected line with $\mathrm{span}(\one)$
gives the convex-hull form.
\end{solution}

\begin{remark}[What this recovers]
For the single white dot $\pm e_j$ of the base paper, $\fdc(\pm e_j)=\pm1/\sqrt d$, so the class gap is $2/\sqrt d$,
equivalently the max-margin is $\gamma_{\mathrm{inv}}=1/\sqrt d$ under the half-gap convention of Theorem~\ref{thm:A}.
The unconstrained margin on that two-image set is $1$, so the ratio is $1/\sqrt d$: invariance can cost a factor
$\sqrt d$ in margin when the signal is purely localized.
\end{remark}

\researchmove{Quotient out the nuisance}{Imposing invariance is the same as projecting the data onto the invariant
subspace and solving the ordinary problem there. The constraint did not make the problem harder; it made the data
lower-dimensional. This ``invariant problem $=$ ordinary problem on the projected data'' move recurs for general
finite groups (replace $\Piinv$ by $P_G$) and is the reason the convex-hull picture keeps working.}

\begin{problem}[General finite group]\label{pr:finite}
Generalize: for a finite orthogonal group $\{R_g\}$ with average projector $P_G$, show the invariant linear max-margin
of $X_\pm$ equals $\tfrac12\dist(\conv P_GX_+,\conv P_GX_-)$.
\end{problem}
\begin{hints}
\item Repeat Problem~\ref{pr:A}(a) with $R_g$ in place of $S_s$: invariance $\Leftrightarrow w\in V^G$.
\item For $w\in V^G$, $w^\top x=w^\top P_Gx$ (since $w=P_Gw$ and $P_G$ is self-adjoint), so the classifier only sees $P_Gx$.
\item Apply Theorem~\ref{thm:hull} to the projected sets.
\end{hints}
\begin{solution}
Invariance of $w^\top x$ for all $x$ forces $R_g^\top w=w$ for all $g$, i.e.\ $w\in V^G$. For such $w$,
$w^\top x=(P_Gw)^\top x=w^\top P_Gx$, so every invariant linear classifier on $X_\pm$ is an ordinary linear classifier
on $P_GX_\pm$ with the same normal and signed margins. Theorem~\ref{thm:hull} on $A=P_GX_+,B=P_GX_-$ gives the result.
The cyclic case is $P_G=\Piinv$.
\end{solution}

\section{\texorpdfstring{The $\eta/L$ certificate, its failure of converse, and the two-sided bracket}{The eta/L certificate, its failure of converse, and the two-sided bracket}}\label{sec:etaL}

The governing scalar of the whole project is $\eta/L$: the margin $\eta$ of an invariant feature divided by its
Lipschitz constant $L$. By Definition~\ref{def:lipmargin} it certifies $r_2\ge\eta/L$. The subtle and important
point is how tight that is.

\begin{problem}[No universal converse]\label{pr:noconv}
Exhibit an invariant feature and a data point where the true robust radius is arbitrarily larger than $\eta/L$, so
no inequality $r_2\le C\,\eta/L$ can hold with a universal constant. (Hint: use a feature that depends only on the
DC coordinate but is very flat near the boundary.)
\end{problem}
\begin{hints}
\item Work on the DC line $t=\fdc(x)$; take the invariant feature $\Psi_n(x)=t^{2n+1}$.
\item Put $x_\pm=\pm e$ with $e=d^{-1/2}\one$, labels $\pm1$. Compute the feature margin $\eta$ and the Lipschitz
constant of $t\mapsto t^{2n+1}$ on $[-1,1]$.
\item The decision boundary is $t=0$, at Euclidean distance $1$ from $e$. Compare to $\eta/L$.
\end{hints}
\begin{solution}
With $g_n=\Psi_n$ and $x_\pm=\pm e$, the feature margin is $\eta=1$. On $[-1,1]$, $|\Psi_n'(t)|=(2n+1)t^{2n}$ peaks at
$L_n=2n+1$, so $\eta/L_n=1/(2n+1)\to0$. But the boundary $t=0$ is at Euclidean distance $\|e-0\|_2=1$ from $e$, and
moving within $\Vac$ does not change $t$ at all, so $r_2(g_n,e)=1$. The ratio $r_2/(\eta/L_n)=2n+1$ is unbounded.
Hence $\eta/L$ is a certificate and a predictor, not a characterization, in general.
\end{solution}

\discover{when a converse could hold}{The counterexample worked because the feature went \emph{flat} near the boundary
(its slope, the lower Lipschitz rate, fell to zero). Before reading on, ask: \emph{what extra quantity would prevent
this?} If the score were guaranteed to change at least at rate $\alpha$ along the path to the boundary, what radius
upper bound would follow? You are about to invent the co-Lipschitz constant.}

\begin{theorem}[Two-sided bracket: a conditional converse]\label{thm:sandwich}
Suppose $g$ is correct at $x$ with $\mu=y\,g(x)>0$, is $L$-Lipschitz toward the boundary, and admits a path $\gamma$
to the boundary along which $|g(\gamma(t))-g(x)|\ge\alpha\|\gamma(t)-x\|$ (a \emph{co-Lipschitz}, or lower-Lipschitz,
bound). Then $\mu/L\le r_2\le\mu/\alpha$ (the score margin $\mu$ here plays the role of the feature margin $\eta$, so
this is the $\eta/L$ certificate completed to a bracket), with condition number $\kappa=L/\alpha\ge1$; the certificate
is exact ($r_2=\mu/L$, $\kappa=1$) for linear invariant features.
\end{theorem}

\begin{problem}[Prove the bracket]\label{pr:sandwich}
Prove both arms, show $\kappa\ge1$, and check the linear case gives $\kappa=1$. Then reconcile with
Problem~\ref{pr:noconv}: which constant blows up there?
\end{problem}
\begin{hints}
\item Lower arm is Definition~\ref{def:lipmargin}.
\item Upper arm: at the first boundary hit $\gamma(T)$, $g=0$, so $\mu=|g(x)-g(\gamma(T))|\ge\alpha\|\gamma(T)-x\|$.
\item The boundary point is a witness, so $r_2\le\|\gamma(T)-x\|\le\mu/\alpha$.
\item Linear $g=w^\top x$: along $\pm w$ the secant slope equals $\|w\|_2$, so $L=\alpha=\|w\|_2$.
\item In Problem~\ref{pr:noconv}, the secant co-Lipschitz is $\alpha=1$ but $L=2n+1$, so $\kappa=2n+1$ is what is unbounded.
\end{hints}
\begin{solution}
Lower arm: no perturbation under $\mu/L$ flips the sign, so $r_2\ge\mu/L$. Upper arm: at the boundary hit
$g(\gamma(T))=0$, so $\mu=|g(x)|=|g(x)-g(\gamma(T))|\ge\alpha\|\gamma(T)-x\|$, giving $\|\gamma(T)-x\|\le\mu/\alpha$;
since $\gamma(T)$ flips the prediction, $r_2\le\mu/\alpha$. The secant slope cannot exceed the Lipschitz bound, so
$\alpha\le L$ and $\kappa\ge1$. For linear $g$ the straight path to the hyperplane has $L=\alpha=\|w\|_2$, so $\kappa=1$
and $r_2=\mu/\|w\|_2$ exactly. In the cubic example the straight DC path has $\alpha=1$ (the feature drops from $1$ to
$0$ over distance $1$) while $L=2n+1$; the bracket $[1/(2n+1),1]$ correctly contains $r_2=1$, at the $\mu/\alpha$ end,
and the failure of a universal converse is exactly the unboundedness of $\kappa$ over the model class.
\end{solution}

\researchmove{Bracket, do not chase equality}{When an inequality has no converse, the move is not to give up but to
name the missing ingredient and bound the gap by it. Here the gap is a condition number $\kappa=L/\alpha$. This both
rescues $\eta/L$ as an approximate characterization and explains the counterexample as the $\kappa\to\infty$ corner.
The same instinct, ``what one constant controls the gap,'' recurs in the lazy-regime theorem, where the gap is the
orbit-variance.}

\section{A convolution-square-pool layer computes the power spectrum}\label{sec:ps}

Linear invariance is weak (Theorem~\ref{thm:A}). The escape is a nonlinear invariant feature an actual network can
compute.

\begin{lemma}[Power-spectrum span]\label{lem:ps}
A single layer of circular convolution, pointwise squaring, and global sum-pooling realizes exactly the
shift-invariant quadratic features, and these span the power-spectrum coordinates $|\hat x_k|^2$. On the ball
$\|x\|_2\le B$ the map $Q(x)=(|\hat x_k|^2)_k$ is $2B$-Lipschitz.
\end{lemma}

\begin{problem}[Prove the span and the Lipschitz bound]\label{pr:ps}
(a) Show $\sum_j (w\star x)_j^2$, with $w\star x$ the circular cross-correlation, is shift-invariant and equals a
weighted sum of $|\hat x_k|^2$. (b) Show $\|Q(x)-Q(y)\|_2\le 2B\|x-y\|_2$ on the ball.
\end{problem}
\begin{hints}
\item (a) Circular correlation is diagonal in Fourier: $\widehat{w\star x}_k=\overline{\hat w_k}\,\hat x_k$. Use Parseval.
\item (a) $\sum_j(w\star x)_j^2=\sum_k|\hat w_k|^2|\hat x_k|^2$, a nonnegative combination of power coordinates; varying $w$ sweeps the coordinates.
\item (b) For complex $z,z'$, $\big||z|^2-|z'|^2\big|\le|z-z'|(|z|+|z'|)$; apply coordinatewise to $\hat x,\hat y$ and use unitarity.
\end{hints}
\begin{solution}
(a) By the convolution theorem and Parseval, $\sum_j(w\star x)_j^2=\sum_k|\widehat{w\star x}_k|^2=\sum_k|\hat w_k|^2|\hat x_k|^2$;
this is shift-invariant (each $|\hat x_k|^2$ is) and is a nonnegative-weighted sum of power coordinates, so a bank of
such units spans the power spectrum. (b) Coordinatewise $\big||\hat x_k|^2-|\hat y_k|^2\big|\le|\hat x_k-\hat y_k|(|\hat x_k|+|\hat y_k|)$;
summing, Cauchy--Schwarz and unitarity of the DFT give $\|Q(x)-Q(y)\|_2\le\|x-y\|_2(\|x\|_2+\|y\|_2)\le2B\|x-y\|_2$.
\end{solution}

\begin{remark}[Phase-blindness and the bispectrum escape]
The power spectrum is invariant but cannot separate signals that differ only in phase, such as a spike from its
reflection (both have flat $|\hat x_k|^2$). The degree-three bispectrum $B_{k_1,k_2}=\hat x_{k_1}\hat x_{k_2}\overline{\hat x_{k_1+k_2}}$
is shift-invariant and phase-retaining, and separates the $\pm e_j$ dot with margin of order $d^{-3/2}$ in a single
bispectrum coordinate (the full bispectrum vector in Euclidean norm scales differently), at the cost of a larger
Lipschitz constant, hence a smaller certified $\eta/L$. This is the ``go to higher degree'' answer to the
find-the-question prompt in \S\ref{sec:group}.
\end{remark}

\section{Orbit-bundle form, orbit rank, and the degeneracy correction}\label{sec:orbit}

Two structural facts make the optimization analysis possible and also explain why the textbook examples mislead.

\begin{lemma}[Orbit-bundle form]\label{lem:bundle}
A convolution-plus-global-pooling channel equals a two-layer network whose hidden units are tied into a shift orbit
with an orbit-averaged readout: $f(x)=\sum_r a_r\frac1d\sum_s\sigma(\langle w_r,S_sx\rangle)$. Imposing exact
shift-invariance on a two-layer ReLU network is exactly this weight-tying.
\end{lemma}

\begin{lemma}[Rank of an orbit]\label{lem:rank}
The cyclic orbit $\{S_su\}_s$ spans a subspace of dimension equal to the number of nonzero Fourier coefficients of
the base $u$.
\end{lemma}

\begin{problem}[Prove the rank formula and find the two degenerate bases]\label{pr:rank}
(a) Prove Lemma~\ref{lem:rank}. (b) Show that the two ``natural'' single-orbit constructions are each degenerate in
opposite ways: a single cosine of frequency $k$, and a localized spike $e_j$.
\end{problem}
\begin{hints}
\item (a) In Fourier coordinates a shift multiplies coordinate $k$ by a phase; the orbit lives in the span of the
frequencies where $\hat u_k\ne0$, and the distinct phase vectors are linearly independent there.
\item (b) A pure cosine has only two nonzero Fourier coefficients (at $\pm k$): orbit rank $2$, and its signed mean is zero.
\item (b) A spike has all coefficients nonzero (full rank, orthonormal orbit) but flat power spectrum, so the
power-spectrum invariant cannot separate two spike orbits.
\end{hints}
\begin{solution}
(a) Write $u=\sum_{k:\hat u_k\ne0}\hat u_k v_k$ with $v_k$ the Fourier modes; $S_su=\sum_k \hat u_k e^{-2\pi isk/d}v_k$.
The vectors $\{S_su\}_s$ therefore lie in $\mathrm{span}\{v_k:\hat u_k\ne0\}$, and the Vandermonde structure of the
phases makes them span it, so the rank is $\#\{k:\hat u_k\ne0\}$. (b) A cosine $\cos(2\pi kj/d)$ has nonzero
coefficients only at $\pm k$: rank $2$, far from the near-orthogonal $\Omega(d)$-cluster geometry the
feature-averaging theorems assume, and its signed orbit mean (the direction those theorems use) is zero. A spike has
all nonzero coefficients, so its orbit is full rank and orthonormal (the cluster assumption holds), but its power
spectrum is flat, so $\Sep(Q)=0$ and the power-spectrum (quadratic) invariant classifier cannot separate two spike classes (a
higher-degree invariant such as the bispectrum still can; see \S\ref{sec:ps}). Neither
base is a valid test of ``does invariance re-point the implicit bias''; generic full-rank, phase-rich bases are.
\end{solution}

\researchmove{Find the degenerate case before you trust the example}{The textbook constructions are exactly the
measure-zero cases where the premise quietly fails. Before importing a theorem (here, feature-averaging
non-robustness), check that your data satisfies its hypotheses; the rank formula is the cheap diagnostic that the
single-orbit examples do not.}

\section{Self-symmetrization: invariance is margin-free on orbit-closed data}\label{sec:selfsym}

Now the first optimization-side theorem. It is a clean, surprising negative result: on orbit-closed data, hard-wiring
shift-invariance does not change the achievable margin at all.

\begin{theorem}[Self-symmetrization]\label{thm:selfsym}
Let the data be orbit-closed with labels constant on orbits. For two-layer ReLU networks
$f(x)=\sum_r a_r\sigma(w_r^\top x)$ (bias-free, degree-2 homogeneous), the minimum parameter norm
$\tfrac12\sum_r(a_r^2+\|w_r\|^2)$ achieving unit margin is the same for the unconstrained class and the
shift-invariant (orbit-tied) class. Imposing exact shift-invariance changes neither the max-margin value nor the
achievable margin.
\end{theorem}

\begin{problem}[Prove self-symmetrization]\label{pr:selfsym}
Two inequalities. (a) The easy direction. (b) The reverse: take any unconstrained solution, symmetrize it, and show
the parameter norm does not grow. Use Problem~\ref{pr:relu}.
\end{problem}
\begin{hints}
\item (a) The tied class is a subset of the unconstrained class, so its min-norm is at least as large.
\item (b) Given an unconstrained $f^\star$ with margin $\ge1$, define $f_{\mathrm{sym}}(x)=\frac1d\sum_s f^\star(S_sx)$. Why is it invariant, in the tied class, and still margin $\ge1$?
\item (b) Margin: orbit-closure means each $S_sx_i$ is a training point of the same label, so the average stays $\ge1$.
\item (b) Norm: $f_{\mathrm{sym}}$ is a bundle of units $\{(a_r/d,\,S_{-s}w_r)\}$; sum the per-unit min-norm $|a_r/d|\,\|w_r\|$ over the $d$ shifts using Problem~\ref{pr:relu}.
\end{hints}
\begin{solution}
(a) Tied $\subseteq$ unconstrained, so $\min$-norm$_{\mathrm{tied}}\ge\min$-norm$_{\mathrm{unc}}$. (b) Let $f^\star$ be
unconstrained with margin $\ge1$. Its orbit average $f_{\mathrm{sym}}(x)=\frac1d\sum_s f^\star(S_sx)$ is
shift-invariant (reindex the pooling sum) and lies in the tied class by Lemma~\ref{lem:bundle}. Orbit-closure with
constant labels gives $y_if^\star(S_sx_i)\ge1$ for all $s$, so $y_if_{\mathrm{sym}}(x_i)\ge1$: margin preserved. By
Problem~\ref{pr:relu} each unit's min-norm is $|a|\,\|w\|$. The bundle realizing $f_{\mathrm{sym}}$ has units
$\{(a_r/d,S_{-s}w_r)\}_{r,s}$ whose total min-norm is $\sum_r\sum_s \frac{|a_r|}{d}\|w_r\|=\sum_r|a_r|\|w_r\|$
(each of the $d$ shifts contributes $\frac{|a_r|}{d}\|w_r\|$, and shifts preserve $\|w_r\|$), equal to the min-norm of
$f^\star$. Hence
$\min$-norm$_{\mathrm{tied}}\le\min$-norm$_{\mathrm{unc}}$, and the two are equal.
\end{solution}

\begin{remark}[What it does and does not give]
Self-symmetrization explains the experimental finding that on generic orbit data the unconstrained and tied networks
reach the same robustness: gradient descent ``self-symmetrizes.'' It does \emph{not} fix which max-margin solution is
selected, and the certified radius $\eta/L$ depends on the Lipschitz term $L$, not only the margin. That selection is
where invariance can still help, on cluster geometry; see \S\ref{sec:lazy}.
\end{remark}

\section{The two-orbit quadratic surrogate has an exact margin}\label{sec:quad}

\begin{theorem}[Two-orbit quadratic surrogate]\label{thm:quad}
For two-orbit data $X_+=\Orb(u)$, $X_-=\Orb(v)$ with distinct power spectra $Q(u)\ne Q(v)$, the shift-invariant
quadratic surrogate has max-margin $\eta_Q=\tfrac12\|Q(u)-Q(v)\|_2$ and, on $\|x\|_2\le B$, certified radius
$\ge\eta_Q/(2B)$.
\end{theorem}

\begin{problem}[Derive the exact margin]\label{pr:quad}
Show that in power-spectrum feature space each class collapses to a single point, and apply convex-hull duality and
Lemma~\ref{lem:ps}.
\end{problem}
\begin{hints}
\item The power spectrum is shift-invariant, so every member of $\Orb(u)$ maps to the same point $Q(u)$.
\item Two single points: the max-margin is half their distance (Theorem~\ref{thm:hull}).
\item The radius certificate is Lemma~\ref{lem:ps}'s $2B$-Lipschitz bound through Definition~\ref{def:lipmargin}.
\end{hints}
\begin{solution}
By shift-invariance $Q(S_su)=Q(u)$, so $X_+$ and $X_-$ map to the singletons $\{Q(u)\}$ and $\{Q(v)\}$. The maximum
margin between two distinct points is half their Euclidean distance, $\eta_Q=\tfrac12\|Q(u)-Q(v)\|_2$, attained by the
bisecting hyperplane. Lemma~\ref{lem:ps} gives the feature map Lipschitz constant $2B$ on the ball, so
Definition~\ref{def:lipmargin} yields radius $\ge\eta_Q/(2B)$. (One checks the optimal head lies in the real-filter
span, so the surrogate is realizable.)
\end{solution}

\section{The lazy-regime cluster theorem}\label{sec:lazy}

This is the positive optimization result: in the lazy regime, on cluster-geometry data, the invariant network reaches
an $\eta/L$ at least as large as the unconstrained one, and strictly larger when the unconstrained interpolant carries
non-invariant energy. The proof is the non-expansiveness of Problem~\ref{pr:proj} transported to the tangent kernel's
function space.

\begin{theorem}[Invariance lowers the Lipschitz constant in the lazy regime]\label{thm:lazy}
In the lazy/NTK regime, with a smooth activation so the tangent feature map $\Phi$ is $L_\Phi$-Lipschitz, on
orbit-closed orbit-labeled data with a $G$-invariant tangent kernel: the orbit-averaging operator $\Pi_G$ is the
orthogonal projection of the tangent RKHS $\mathcal H_K$ onto its invariant subspace, hence non-expansive; the
invariant min-norm interpolant satisfies $\|f_{\mathrm{inv}}\|_K\le\|f_{\mathrm{unc}}\|_K$ at equal margin; therefore
its RKHS-norm-based Lipschitz bound, hence the certified $\eta/L$, is no worse, and strictly better when the
RKHS-norm-squared gap $\|(I-\Pi_G)f_{\mathrm{unc}}\|_K^2$ (the orbit-variance) is positive. This gap is the RKHS-norm
gap; it bounds, but need not equal, the exact robust-radius or $\eta/L$ gap.
\end{theorem}

\begin{problem}[Assemble the lazy-regime proof]\label{pr:lazy}
You are given: (i) lazy $\Rightarrow$ the trained net is the min-RKHS-norm interpolant of a fixed kernel; (ii) the RKHS
Lipschitz bound $|f(x)-f(y)|\le\|f\|_K\|\Phi(x)-\Phi(y)\|$. Build the theorem in three steps, then state the one place
the smooth-activation hypothesis is load-bearing.
\end{problem}
\begin{hints}
\item Step 1: $\Pi_G$ is self-adjoint and idempotent on $\mathcal H_K$ (it is the Reynolds projector, Problem~\ref{pr:proj}), so it is an orthogonal projection: $\|f\|_K^2=\|\Pi_Gf\|_K^2+\|(I-\Pi_G)f\|_K^2$.
\item Step 2: orbit-constant labels make $\Pi_Gf_{\mathrm{unc}}$ still interpolate with the same margin, so it is feasible for the invariant min-norm problem; hence $\|f_{\mathrm{inv}}\|_K\le\|\Pi_Gf_{\mathrm{unc}}\|_K\le\|f_{\mathrm{unc}}\|_K$.
\item Step 3: smaller RKHS norm gives smaller Lipschitz constant via (ii) and the $L_\Phi$-Lipschitz feature map; at equal margin $\eta/L$ improves.
\item Load-bearing: step 3 needs $\Phi$ Lipschitz. The bare two-layer ReLU NTK feature map is \emph{not} Lipschitz (its RKHS has unit-norm functions of unbounded Lipschitz constant), so the clean statement needs a smooth activation.
\end{hints}
\begin{solution}
Step 1: a $G$-invariant kernel makes $\Pi_G$ self-adjoint and idempotent on $\mathcal H_K$ with spectrum $\{0,1\}$,
hence the orthogonal projection onto the invariant subspace, with the Pythagorean identity. Step 2: orbit-constancy
gives $y_if_{\mathrm{unc}}(g\cdot x_i)\ge1$ for all $g$, so $y_i(\Pi_Gf_{\mathrm{unc}})(x_i)\ge1$; thus
$\Pi_Gf_{\mathrm{unc}}$ is feasible for the invariant min-norm problem, giving $\|f_{\mathrm{inv}}\|_K\le\|\Pi_Gf_{\mathrm{unc}}\|_K$,
which Pythagoras bounds by $\|f_{\mathrm{unc}}\|_K$. Step 3: the RKHS Lipschitz bound with an $L_\Phi$-Lipschitz $\Phi$
gives $\mathrm{Lip}(f)\le L_\Phi\|f\|_K$; substituting the norm inequality, $L_{\mathrm{inv}}\le L_{\mathrm{unc}}$ at
equal margin, so $\eta/L$ does not decrease, and the gap is the discarded non-invariant energy
$\|(I-\Pi_G)f_{\mathrm{unc}}\|_K^2$, the orbit-variance. The smoothness is needed only in step 3; for bare ReLU the
feature map is non-Lipschitz, so one restricts to smooth activations or carries $\sup\|\nabla\Phi\|$ explicitly. The
rich-regime ReLU trajectory remains open.
\end{solution}

\researchmove{Transport a finite-dimensional lemma to function space}{The whole theorem is Problem~\ref{pr:proj}(iii)
``a projection cannot lengthen a vector'' moved from $\R^d$ to the RKHS, plus the dictionary norm $\to$ Lipschitz. When
a clean linear-algebra fact stalls because the objects are now functions, look for the Hilbert space where the same
words are true. The honest scoping (smooth activation, lazy regime) is part of the result, not an apology.}

\section{The open problem: which solution does rich-regime training select?}\label{sec:open}

Self-symmetrization (\S\ref{sec:selfsym}) says invariance is margin-free on generic orbit data; the lazy-regime
theorem (\S\ref{sec:lazy}) says it lowers $L$ on cluster geometry. The gap between them is the live research question.

\begin{conjecture}[Trajectory-level selection]\label{conj:opt}
On nondegenerate cluster-geometry orbit data there is a regime (initialization scale, width) where rich-regime
gradient flow on the shift-invariant network reaches a KKT point of strictly larger $\eta/L$ than the unconstrained
network, even though both attain perfect clean margin.
\end{conjecture}

The checklist to attack it (the project's ``what to prove''):
\begin{itemize}[leftmargin=1.4em,nosep]
\item[(P1)] Write the invariant max-margin (KKT) program on two-orbit data and characterize the selected $\eta/L$.
\item[(P2)] Show the unconstrained gradient flow reaches a strictly smaller $\eta/L$ on the same data, \emph{or}
exhibit the minimal orbit model where the feature-averaging premise provably holds.
\item[(P3)] Combine into Conjecture~\ref{conj:opt}: weight-sharing re-points the bias toward larger $\eta/L$.
\item[(P4)] Decide whether the non-robust-selection premise holds on \emph{any} full-rank orbit model at all.
\end{itemize}

\discover{the right minimal model}{The experiments show invariance helps on ``near-orthogonal cluster geometry'' but
not on generic single orbits. Before reading the project notes, ask: \emph{what is the smallest data model that is
simultaneously (a) full-rank/phase-rich enough for the invariant feature to separate the classes and (b) cluster-like
enough for the feature-averaging non-robustness to bite?} The two requirements pull in opposite directions; finding
the model where both hold is most of the open problem.}

% =====================================================================
\part*{Part III: How to find the questions yourself}
\addcontentsline{toc}{section}{Part III: How to find the questions yourself}

Parts I and II trained you to follow and reconstruct. This part trains the harder skill: generating the question.
We do it by reverse-engineering how each Part~II result was actually found, then giving you open prompts with no
stated question.

\section{The anatomy of the project's results}

Every result in Part~II came from one of a small number of moves. Naming them lets you reuse them.

\begin{enumerate}[leftmargin=1.6em]
\item \textbf{Quotient out the nuisance} (\S\ref{sec:thmA}). The question ``how much margin survives invariance?'' became
``solve the ordinary problem on the projected data.'' Whenever a constraint is a symmetry, project and re-ask.
\item \textbf{Diagonalize the symmetry} (\S\ref{sec:group}, \S\ref{sec:ps}). Hard statements about shifts became one-liners
in Fourier coordinates. When one group acts, change to its eigenbasis before doing anything else.
\item \textbf{Bracket what has no converse} (\S\ref{sec:etaL}). A one-sided certificate with no converse was rescued by
naming the missing constant ($\alpha$) and bounding the gap by $\kappa=L/\alpha$.
\item \textbf{Find the degenerate case} (\S\ref{sec:orbit}). The textbook examples were exactly where the premise failed;
the rank formula is the diagnostic that exposes it.
\item \textbf{Reduce to the per-unit atom} (\S\ref{sec:selfsym}). A theorem about whole networks was the one-unit
homogeneity identity (Problem~\ref{pr:relu}) summed over an orbit.
\item \textbf{Transport a finite lemma to function space} (\S\ref{sec:lazy}). ``A projection cannot lengthen a vector''
became the entire lazy-regime theorem once moved to the RKHS.
\end{enumerate}

\section{Two modes of solving, and why the second is rare}

There is a real difference between \emph{following} a decomposition someone hands you and \emph{producing} it. The
first is pattern-completion; the second needs you to (i) notice that a question is even there, (ii) guess the right
object to introduce, and (iii) test it on the smallest example before investing in a proof. The hint ladders in
Part~II quietly performed (ii) and (iii) for you. The exercises below withhold them.

\begin{problem}[Generate the question, then answer it: a new group]\label{pr:newgroup}
We did everything for cyclic shifts. Consider instead the group of \emph{horizontal and vertical} circular shifts on a
$d_1\times d_2$ image, $\Z_{d_1}\times\Z_{d_2}$. No sub-question is stated. Decide what the right analogues of the DC
line, the orbit rank, and the power spectrum are, state the version of Theorem~\ref{thm:A} you expect, and prove the
one you are most confident about.
\end{problem}
\begin{hints}
\item Which move applies first? (Diagonalize: the 2D DFT diagonalizes both shift directions at once.)
\item The fixed subspace of all 2D shifts is still the constants; the orbit rank is the number of nonzero 2D Fourier
coefficients; the power spectrum is $|\hat x_{k_1,k_2}|^2$.
\item Theorem~\ref{thm:A} goes through verbatim with $\Piinv$ the 2D DC projector; the proof is unchanged because it
only used that $P_G$ projects onto the constants.
\end{hints}

\begin{problem}[Generate the question: a different threat]\label{pr:newthreat}
The bracket theorem (\S\ref{sec:etaL}) was stated for $\ell_2$. Without being told what to prove, work out what changes
under an $\ell_\infty$ threat, what the right co-Lipschitz constant is, and whether the linear case is still exact.
State a conjecture and test it on $g(x)=w^\top x$.
\end{problem}
\begin{hints}
\item The certificate uses the dual norm (Problem~\ref{pr:dual}); for $\ell_\infty$ that is the $\ell_1$ gradient.
\item The linear exact case becomes $r_\infty=\mu/\|w\|_1$; check this is the point-to-hyperplane distance in $\ell_\infty$.
\end{hints}

\begin{problem}[Generate the question: where could invariance \emph{help} robustness more?]\label{pr:help}
Self-symmetrization says invariance is margin-free on generic orbit data, and the lazy theorem says it helps via $L$
on cluster geometry. Pose the sharpest experiment you can that would distinguish ``invariance helps through margin''
from ``through Lipschitz,'' design its arms, and predict the outcome from the theory before any run.
\end{problem}

\section{Taste: which questions are worth asking}

Not every well-posed question is worth a month. The ones that paid off here shared three features: they had a
\emph{small testable instance} (a $d=4$ orbit, a two-image set), they sat at a \emph{contradiction in the literature}
(invariance reported as both helping and hurting robustness), and their answer \emph{changed a decision} (which
architecture to trust, which certificate to report). When you generate questions in Part~III, score each on those
three before committing. A question with no small instance is usually premature; a question whose answer changes
nothing is usually not worth the proof.

% =====================================================================
\part*{Part IV: Experiments and where to contribute}
\addcontentsline{toc}{section}{Part IV: Experiments and where to contribute}

\section{The experimental backbone (the part of the project you can run)}

The theory is tested by a capacity-matched architectural dissection. Four architectures share a parameter count and
differ only in the shift-invariance operator: a \emph{standard} network, an \emph{anti-aliased} network (a fixed
low-pass filter before each subsampling), an \emph{exactly invariant} network (adaptive polyphase sampling, invariant
to all integer shifts), and a \emph{shift-augmented} network (invariance learned from data, not built in). The
measured quantities are the ones from Part~II: shift-consistency, the $\ell_2$ robust radius (a minimum-norm attack),
and the $\eta/L$ decomposition into margin and Lipschitz terms.

The headline empirical facts, which you should be able to predict from the theory:
\begin{itemize}[leftmargin=1.4em,nosep]
\item Under standard training, $\eta/L$ predicts the robust radius while shift-consistency does not.
\item Under adversarial training, shift-consistency \emph{anti}-predicts robustness, and the threat-matched ratio
$\eta/\|\nabla M\|_1$ predicts it; the exactly invariant arm is the least robust, through reduced margin.
\end{itemize}
A new collaborator's natural first experimental task is to reproduce one cell of this dissection and read its results
file, then re-derive which Part~II result each number tests.

\section{Open problems ranked by accessibility}
\begin{enumerate}[leftmargin=1.6em]
\item \textbf{(P1) Two-orbit invariant KKT in closed form} (\S\ref{sec:open}). Self-contained, frequency-domain,
a good first theory contribution. Start from Theorem~\ref{thm:quad} and replace the quadratic surrogate by ReLU.
\item \textbf{The minimal cluster-orbit model} (find-the-question, \S\ref{sec:open}). Decide whether the
feature-averaging premise can hold on any full-rank orbit model; if not, the conjecture must be restricted.
\item \textbf{Rich-regime selection} (Conjecture~\ref{conj:opt}). The hardest: track which KKT point gradient flow
selects when several exist. Needs trajectory information, not just the variational limit.
\item \textbf{$\ell_\infty$ bracket and tighter $\kappa$ estimates} (Problem~\ref{pr:newthreat}). Extend the
two-sided certificate and estimate $\kappa$ for the power-spectrum model.
\end{enumerate}

\section{Reading roadmap}
\begin{enumerate}[leftmargin=1.6em,nosep]
\item This primer, Parts I--II, with the problems done closed-book.
\item The paper \texttt{main.tex} end to end; you should recognize every theorem from Part~II.
\item The technical report \texttt{theory\_full\_checked\_v2.tex} for the reference proofs and the optimization
roadmap section.
\item One experimental cell reproduced; one open problem from Part~IV attempted for two weeks before proposing.
\end{enumerate}

\section*{Glossary}
\addcontentsline{toc}{section}{Glossary}
\begin{description}[leftmargin=1.2em,style=nextline]
\item[Robust radius $r_p(f,x)$] distance from $x$ to the nearest point of opposite prediction, in $\ell_p$.
\item[Certificate] a guaranteed lower bound on the robust radius (versus an attack, an upper bound).
\item[$\eta/L$] invariant-feature margin divided by its Lipschitz constant; certifies $r_2\ge\eta/L$.
\item[Co-Lipschitz constant $\alpha$] a guaranteed lower rate of change toward the boundary; gives the converse arm.
\item[Condition number $\kappa=L/\alpha$] controls how far $\eta/L$ is from the true radius; $\kappa=1$ is exact.
\item[Orbit] the set of all shifts of a signal; orbit-closed means closed under all shifts.
\item[DC / AC subspace] the constant line $\mathrm{span}(\one)$ and its zero-sum complement.
\item[Power spectrum / bispectrum] degree-two phase-blind, and degree-three phase-retaining, shift invariants.
\item[Reynolds projector $P_G$] the group average $|G|^{-1}\sum_g R_g$; the orthogonal projector onto the invariants.
\item[Self-symmetrization] on orbit-closed data, gradient descent reaches an invariant max-margin solution on its own.
\item[Lazy / rich regime] very-wide fixed-kernel training versus small-init feature-learning training.
\item[Feature averaging] the implicit-bias mechanism by which max-margin solutions become non-robust on cluster data.
\end{description}

\section*{References}
\addcontentsline{toc}{section}{References}
The primer names these works in prose; full details below. The project paper and technical report carry the complete
bibliography.
\sloppy
\begin{description}[leftmargin=1.2em,style=nextline]
\item[Ge, Singla, Basri, Jacobs (2021)] Shift invariance can reduce adversarial robustness. arXiv:2103.02695. (The base paper.)
\item[Kamath, Deshpande, Subrahmanyam, Balasubramanian (2021)] Can we have it all? On the trade-off between spatial and adversarial robustness. arXiv:2002.11318.
\item[Tram\`er, Behrmann, Carlini, Papernot, Jacobsen (2020)] Fundamental tradeoffs between invariance and sensitivity to adversarial perturbations. arXiv:2002.04599.
\item[Zhang (2019)] Making convolutional networks shift-invariant again (anti-aliasing / BlurPool). arXiv:1904.11486.
\item[Chaman, Dokmani\'c (2021)] Truly shift-invariant convolutional neural networks (adaptive polyphase sampling). arXiv:2011.14214.
\item[Soudry, Hoffer, Nacson, Gunasekar, Srebro (2018)] The implicit bias of gradient descent on separable data. arXiv:1710.10345.
\item[Lyu, Li (2020)] Gradient descent maximizes the margin of homogeneous neural networks. arXiv:1906.05890.
\item[Ji, Telgarsky (2020)] Directional convergence and alignment in deep learning. arXiv:2006.06657.
\item[Jacot, Gabriel, Hongler (2018)] Neural tangent kernel: convergence and generalization in neural networks. arXiv:1806.07572.
\item[Chizat, Oyallon, Bach (2019)] On lazy training in differentiable programming. arXiv:1812.07956.
\item[Frei, Vardi, Bartlett, Srebro (2023)] The double-edged sword of implicit bias: generalization vs.\ robustness in ReLU networks. arXiv:2303.01456.
\item[Li, Pan, Lyu, Li (2024)] Feature averaging: an implicit bias of gradient descent leading to non-robustness. arXiv:2410.10322.
\item[Elesedy (2021)] Provably strict generalisation benefit for invariance in kernel methods. arXiv:2106.02346.
\item[Bietti, Mairal (2019)] On the inductive bias of neural tangent kernels. arXiv:1905.12173.
\item[Hein, Andriushchenko (2017)] Formal guarantees on the robustness of a classifier against adversarial manipulation. arXiv:1705.08475.
\item[Tsuzuku, Sato, Sugiyama (2018)] Lipschitz-margin training. arXiv:1802.04034.
\item[Croce, Hein (2020)] Reliable evaluation of adversarial robustness with an ensemble of diverse parameter-free attacks (AutoAttack). arXiv:2003.01690.
\item[Rony et al.\ (2019)] Decoupling direction and norm for efficient gradient-based $\ell_2$ adversarial attacks (DDN). arXiv:1811.09600.
\item[Madry, Makelov, Schmidt, Tsipras, Vladu (2018)] Towards deep learning models resistant to adversarial attacks (PGD adversarial training). arXiv:1706.06083.
\item[Rice, Wong, Kolter (2020)] Overfitting in adversarially robust deep learning. arXiv:2002.11569.
\item[Kakarala (2012)] The bispectrum as a source of phase-sensitive invariants for Fourier descriptors. arXiv:0902.0196.
\end{description}
\fussy

\section*{Colophon}
\addcontentsline{toc}{section}{Colophon}
This primer is a teaching companion to the paper and the technical report; all theorem statements and proofs are the
verified versions in \texttt{theory\_full\_checked\_v2.tex}. It follows the project's house writing rules. Corrections
and better worked examples are welcome and should be checked against the technical report before they are added.

\end{document}
